\documentclass[11pt]{article}
\usepackage{amsmath,amssymb,amsthm}
\usepackage[margin=1.15in]{geometry}
\usepackage{hyperref}

\newtheorem{theorem}{Theorem}[section]
\newtheorem{proposition}[theorem]{Proposition}
\newtheorem{lemma}[theorem]{Lemma}
\newtheorem{corollary}[theorem]{Corollary}
\newtheorem{conjecture}[theorem]{Conjecture}
\theoremstyle{definition}
\newtheorem{definition}[theorem]{Definition}
\newtheorem{question}[theorem]{Question}
\newtheorem{remark}[theorem]{Remark}

\newcommand{\N}{\mathbb{N}}
\newcommand{\R}{\mathbb{R}}
\newcommand{\HH}{\mathcal{H}}
\newcommand{\PP}{\mathcal{P}}
\newcommand{\U}{\mathcal{U}}
\newcommand{\D}{\mathcal{D}}
\newcommand{\lk}{\ell^2(\kappa)}
\newcommand{\cod}{\operatorname{codim}}
\newcommand{\supp}{\operatorname{supp}}

\title{Exact paving for Fubini products of normal measures,\\
with an application to maximal projection filters\thanks{The theorems of this paper are machine-checked. The verification artifact is archived at \href{https://doi.org/10.5281/zenodo.21987889}{\texttt{10.5281/zenodo.21987889}} (version~1.0); see \cite{plufLean} for the concept DOI, which resolves to the current version.}}
\author{David V.\ Feldman\thanks{University of New Hampshire;
\texttt{dvfinnh@gmail.com}.}
\and Alexander Wilce\thanks{Susquehanna University.}}
\date{}

\begin{document}
\maketitle

\begin{abstract}
Let $\kappa$ be measurable and $\D$ a normal measure on $\kappa$. We prove
that the Fubini product $\D\otimes\D$ satisfies \emph{exact paving}: for
every assignment of a countable set $c_x$ to each point $x$, some set of
$\D\otimes\D$ mutually avoids the assignment. The proof uses only Rowbottom
homogeneity for finitely many colorings of triples and quadruples; normality
of the product is neither available nor needed. Consequently the block
filter of $\D\otimes\D$ in the projection lattice of $\ell^2(\kappa\times
\kappa)$ is a $\kappa$-complete nonprincipal maximal filter with a singleton
state face --- although $\D\otimes\D$ is not isomorphic to any normal
measure. This answers, negatively, the question of the third paper in this
series: maximality of the block filter does not characterize normality up to
isomorphism. We organize the invariants that do and do not survive: the
full-selector property of normal measures, which holds for arbitrary
partitions, is isomorphism-invariant, and fails for the product; the
$\sigma$-Q-point property, which the product enjoys, by an
injectivity-on-a-square theorem; and an exact-paving property lying between
them.
In a final section we return to $\aleph_0$ and the rigidity conjecture of
the second paper, proving three structural constraints on any addable
blocker of the witness subspace: it meets the block of any $m$ pieces in
dimension at most $m$, its total support belongs to the ultrafilter, and on
every measure-one block it contains vectors of infinite piece-spread.
\end{abstract}

\section{Introduction}\label{sec:intro}

This paper continues a series \cite{FWI,FWII,FWIII} on maximal filters in
the projection lattice $\PP(H)$ of a real Hilbert space (\emph{plufs}), the
lattice-theoretic linearization of ``bounded sequences converge along
ultrafilters.'' The third paper established that at a measurable cardinal
$\kappa$, along a normal measure $\D$, the theory collapses to an exact
form: every bounded operator is exactly diagonal on a measure-one
coordinate block (\emph{exact paving}), every closed subspace either
contains a measure-one block or meets one trivially (\emph{exact
dichotomy}), and consequently the block filter
\[
\Phi(\U)\;=\;\{M\in\PP(\lk):\HH_S\subseteq M\text{ for some }S\in\U\}
\]
is a $\kappa$-complete nonprincipal pluf with a singleton state face. It
also proved a necessity result: maximality of $\Phi(\U)$ forces $\U$ to be a
\emph{$\sigma$-Q-point} (every partition of $\kappa$ into countable pieces
admits a partial selector in $\U$), a property enjoyed by every normal
measure, and asked (\cite{FWIII}, Question 5.5) whether maximality
characterizes the normal measures up to isomorphism, proposing the Fubini
product $\D\otimes\D$ as the test case.

The present paper settles the test case, and with it the question. The
product is not isomorphic to any normal measure
(Proposition~\ref{prop:notiso}), yet its block filter is maximal
(Theorem~\ref{thm:main}). The engine is Theorem~\ref{thm:paving-product}:
$\D\otimes\D$ satisfies exact paving outright. Normality entered the paving
proof of \cite{FWIII} through the diagonal intersection; the product admits
no diagonal intersections, but it inherits from $\D$ something just as
serviceable, Rowbottom homogeneity for finite colorings of finite subsets of
$\kappa$, and homogeneity suffices: any assignment of countable
avoidance-sets, examined through the finitely many \emph{membership
patterns} it can exhibit on triples and quadruples, is forced by
countability to avoid a homogeneous square.

The same method yields an injectivity theorem
(Theorem~\ref{thm:inj}) --- every countable-to-one function on pairs from
$\kappa$ is injective on the pairs from a $\D$-large set --- which gives the
$\sigma$-Q-point property of the product directly.

Since the mechanized development accompanying \cite{FWIII} consumed
normality only through the paving conclusion, the entire positive theory
transfers to the product along Theorem~\ref{thm:paving-product}; we spell
this out in Section~\ref{sec:transfer}, introduce the \emph{exact-paving
property} (EPP) as the invariant the transfer isolates, and record the
resulting picture:
\[
\text{normal}\;\Longrightarrow\;\text{EPP}\;\Longrightarrow\;
\text{$\Phi$ maximal}\;\Longrightarrow\;\text{$\sigma$-Q-point},
\]
so that maximality holds strictly beyond the normal measures. What
\emph{does} separate normal measures from products is a selector property:
normal measures admit full selectors for \emph{arbitrary} partitions with no
piece in the measure (Lemma~\ref{lem:full-selector}), the piece-minima
argument of \cite{FWIII} requiring no countability. The property is
isomorphism-invariant, and the column partition refutes it for the product.
The invariants thus separate: $\D\otimes\D$ is maximal and $\sigma$-Q but
not Q-for-all-partitions.

Section~\ref{sec:rigidity} returns to $\aleph_0$ and the rigidity
conjecture of \cite{FWII}: for a non-selective ultrafilter $\U$ on $\N$
with witness subspace $W$, does $W$ belong to every pluf extension of the
block filter? We prove three constraints on any prospective \emph{addable
blocker} $R$ (a closed subspace with $R\cap W=0$ meeting every member of
the block filter in infinite dimension): a generalized thinness bound
(Lemma~\ref{lem:genthin}), a support lemma
(Lemma~\ref{lem:support}), and a Baire-category piece-spread lemma
(Lemma~\ref{lem:spread}). The conjecture remains open; the constraints
close every finitely-generated route to a counterexample. The measurable-cardinal results of this paper
frame the conjecture from above: along $\D$ or $\D\otimes\D$ there are no
witness subspaces at all, so rigidity, if real, is a phenomenon of small
cardinality.

Throughout, $\kappa$ is a measurable cardinal, $\D$ a normal measure
(nonprincipal $\kappa$-complete normal ultrafilter) on $\kappa$, and
$\D\otimes\D$ the Fubini product on $\kappa\times\kappa$:
$X\in\D\otimes\D$ iff $\{\gamma:\{\beta:(\gamma,\beta)\in X\}\in\D\}\in\D$.
We write $[A]^2$ for the two-element subsets of $A$, identified when
convenient with the ordered pairs $(\gamma,\beta)$, $\gamma<\beta$; since
each section $\{\beta:\beta>\gamma\}$ lies in $\D$, the product concentrates
on the upper triangle, and we pass between the triangle and $[\kappa]^2$
without comment. The Hilbert-space background and the definitions of pluf,
block, block filter and witness subspace are as in \cite{FWI,FWII,FWIII}.

\section{Homogeneity preliminaries}\label{sec:prelim}

We use one classical input.

\begin{theorem}[Rowbottom; see {\cite[Theorem 7.17]{Kanamori}}]
\label{thm:rowbottom}
Let $\D$ be a normal measure on $\kappa$, $n<\omega$, $\lambda<\kappa$, and
$F:[\kappa]^n\to\lambda$. Then there is $H\in\D$ with $F$ constant on
$[H]^n$.
\end{theorem}

Applying Theorem~\ref{thm:rowbottom} to the finite product of finitely many
colorings homogenizes them simultaneously; we use this silently. The
following one-line observation supplies all the ``high pivots'' below: if
$H\in\D$ then $H$ has order type $\kappa$, so for every $\lambda<\kappa$
some (indeed, a tail of) $\beta\in H$ satisfies
$|H\cap\beta|\ge\lambda$.

The selector lemma of \cite{FWIII} holds for arbitrary partitions:
countability of the pieces plays no role in its proof.

\begin{lemma}[Full selectors along normal measures]\label{lem:full-selector}
Let $\D$ be normal and $\kappa=\bigsqcup_{i\in I}\sigma_i$ a partition into
nonempty pieces with $\sigma_i\notin\D$ for every $i$. Then the set
$M=\{\min\sigma_i:i\in I\}$ of piece-minima lies in $\D$; in particular
$\D$ contains a set meeting every piece exactly once.
\end{lemma}

\begin{proof}
If $M\notin\D$ then on $\kappa\setminus M\in\D$ the map
$g(\alpha)=\min(\text{piece of }\alpha)$ is regressive, hence constant on a
set of $\D$ by normality; that set lies inside a single piece,
contradicting $\sigma_i\notin\D$. So $M\in\D$, and $M$ meets each piece
exactly once by disjointness of the pieces and distinctness of their
minima.
\end{proof}

The property in the conclusion --- every partition with no piece in the
ultrafilter admits a selector in the ultrafilter --- mentions no order and
is therefore invariant under isomorphism of ultrafilters (transport along a
bijection of the underlying sets). We refer to it as the
\emph{full-selector property}. The countable-piece specialization is the
$\sigma$-Q-point property of \cite{FWIII}.

\begin{proposition}\label{prop:notiso}
$\D\otimes\D$ fails the full-selector property; consequently
$\D\otimes\D$ is not isomorphic to any normal measure.
\end{proposition}

\begin{proof}
Partition $\kappa\times\kappa$ into the columns
$C_\gamma=\{\gamma\}\times\kappa$. No column lies in $\D\otimes\D$ (the set
of first coordinates with $\D$-large section is a singleton, not in $\D$).
A selector meets each column at most once, so its sections are singletons,
none in $\D$; hence no selector lies in $\D\otimes\D$. Since normal
measures have the full-selector property (Lemma~\ref{lem:full-selector})
and the property is isomorphism-invariant, no isomorph of a normal measure
can be $\D\otimes\D$.
\end{proof}

\section{Injectivity on a square}\label{sec:inj}

\begin{theorem}\label{thm:inj}
Let $\D$ be a normal measure on $\kappa$, $\lambda<\kappa$, and let
$g:[\kappa]^2\to V$ be a function all of whose fibers have size less than
$\max(\lambda,\omega_1)$. Then there is $H\in\D$ with $g$ injective on
$[H]^2$.
\end{theorem}

\begin{proof}
Color the three-element subsets $\{h_1<h_2<h_3\}$ of $\kappa$ by the
equality pattern of $g$ on their three pairs, and the four-element subsets
$\{h_1<h_2<h_3<h_4\}$ by the equality pattern of $g$ on their six pairs;
each coloring has finitely many values. By
Theorem~\ref{thm:rowbottom}, fix $H\in\D$ homogeneous for both. Write
$\mu=\max(\lambda,\omega_1)$; all fibers of $g$ have size $<\mu$, and by
the order-type observation of Section~\ref{sec:prelim} we may pick elements
of $H$ with at least $\mu$ predecessors in $H$ whenever needed.

We claim the homogeneous patterns identify no two pairs; injectivity on
$[H]^2$ follows, since two distinct pairs from $H$ span a three- or
four-element subset of $H$. Suppose some identification held
homogeneously; we derive in each case a fiber of $g$ of size $\ge\mu$.

\emph{Three-element patterns.} If $g(h_1,h_2)=g(h_1,h_3)$ homogeneously,
fix $h_1<h_2$ in $H$ and vary $h_3$ over the tail of $H$ above $h_2$: the
fiber of $g(h_1,h_2)$ contains $\kappa$ pairs. If
$g(h_1,h_2)=g(h_2,h_3)$, fix $h_1<h_2$ and vary $h_3$: again $\kappa$
pairs. If $g(h_1,h_3)=g(h_2,h_3)$, choose $h_2<h_3$ in $H$ with
$|H\cap h_2|\ge\mu$ and vary $h_1$ over $H\cap h_2$: the fiber of
$g(h_2,h_3)$ contains $\mu$ pairs.

\emph{Four-element patterns} (identifications between the two
coordinate-disjoint pairs; identifications sharing a coordinate reduce to
three-element patterns). If $g(h_1,h_2)=g(h_3,h_4)$, fix $h_1<h_2$ and
vary $h_3<h_4$ above $h_2$: $\kappa$ pairs in one fiber. If
$g(h_1,h_3)=g(h_2,h_4)$, fix $h_1<h_2<h_3$ and vary $h_4$: $\kappa$ pairs
share the value $g(h_1,h_3)$. If $g(h_1,h_4)=g(h_2,h_3)$, fix
$h_2<h_3$ and $h_1<h_2$, and vary $h_4$: $\kappa$ pairs share the value
$g(h_2,h_3)$.

Each case contradicts the fiber bound; so all patterns are
``all-distinct,'' and $g$ is injective on $[H]^2$.
\end{proof}

\begin{corollary}\label{cor:product-sigmaQ}
$\D\otimes\D$ is a $\sigma$-Q-point; indeed, for every $\lambda<\kappa$,
every partition of $\kappa\times\kappa$ into pieces of size $<\lambda$
admits a partial selector in $\D\otimes\D$, of the form
$\{(\gamma,\beta):\gamma<\beta\in H\}$ with $H\in\D$.
\end{corollary}

\begin{proof}
The product concentrates on the upper triangle, which we identify with
$[\kappa]^2$. Let $g$ assign to each point its piece; the fibers are the
pieces, of size $<\max(\lambda,\omega_1)<\kappa$. Theorem~\ref{thm:inj}
gives $H\in\D$ with $g$ injective on $[H]^2$; the upper triangle
$S=\{(\gamma,\beta):\gamma,\beta\in H,\ \gamma<\beta\}$ lies in
$\D\otimes\D$ (its section at $\gamma\in H$ is a tail of $H$), and
distinct points of $S$ are distinct pairs of $[H]^2$, hence lie in
distinct pieces.
\end{proof}

Together with Proposition~\ref{prop:notiso}, the corollary separates two
invariants: the $\sigma$-Q-point property does not imply the full-selector
property.

\section{Exact paving for the product}\label{sec:paving}

\begin{theorem}\label{thm:paving-product}
Let $\D$ be a normal measure on $\kappa$ and let each point
$x\in\kappa\times\kappa$ carry a countable set
$c_x\subseteq\kappa\times\kappa$. Then there is $S\in\D\otimes\D$ such
that $y\notin c_x$ for all distinct $x,y\in S$. More generally the
conclusion holds with ``countable'' replaced by ``of size $<\lambda$'' for
any $\lambda<\kappa$.
\end{theorem}

\begin{proof}
Write $\mu=\max(\lambda,\omega_1)$ in the general form; in the countable
case $\mu=\omega_1$. As before, work on the upper triangle, identifying
points with pairs. For each ordered pair of \emph{placement patterns} on a
four-element set $\{h_1<h_2<h_3<h_4\}$ --- a choice of two disjoint pairs
from the four elements, one designated $x$ and the other $y$ --- color the
four-element sets of $\kappa$ by the truth value of ``$y\in c_x$''; there
are finitely many placement patterns. Similarly, for each ordered pair of
distinct pairs from a three-element set $\{h_1<h_2<h_3\}$ (patterns in
which $x$ and $y$ share one coordinate), color the three-element sets by
the truth value of ``$y\in c_x$.'' By Theorem~\ref{thm:rowbottom} fix
$H\in\D$ homogeneous for all these finitely many colorings
simultaneously, and let $S$ be the upper triangle of $H$; as in
Corollary~\ref{cor:product-sigmaQ}, $S\in\D\otimes\D$.

Distinct points $x,y\in S$ span a three- or four-element subset of $H$ in
one of the enumerated patterns, so it suffices to show every pattern is
homogeneously colored ``$y\notin c_x$.'' Suppose some pattern were
homogeneously ``yes''; in each case we produce a single $x$ with
$|c_x|\ge\mu$, contradicting the size bound. The scheme is uniform:
holding $x$'s coordinates fixed at suitably chosen elements of $H$, the
coordinates of $y$ still range over a large set, and homogeneity puts
every resulting $y$ into the one set $c_x$. Each unordered placement must be treated in \emph{both} role assignments,
giving twelve patterns in all. In four of them the large family lives in the
region \emph{between} two fixed coordinates: there one passes first to the
tail of $H$ above the lower coordinate --- a set of $\D$ --- and takes the
pivot inside that tail, so that uncountably many elements of $H$ lie
strictly between.

\emph{Patterns in which $y$ owns the top position $h_4$}, with $x$ on two
of the remaining positions: fix $x$'s coordinates (choosing, when the
pattern requires an element of $H$ strictly between them or below them,
elements of $H$ making those regions nonempty --- possible since $H$ is
unbounded of order type $\kappa$), and vary $h_4$ over a tail of $H$:
$\kappa$-many values of $y$, all in $c_x$.

\emph{Patterns in which $x$ owns the top position and $y$ lies among the
lower positions}: choose $x$'s lower coordinate to be an element $h$ of
$H$ with $|H\cap h|\ge\mu$ (order-type observation), and its top
coordinate above it; the free coordinates of $y$ then range over subsets
of $H\cap h$ (together, when the pattern permits, with a region between
$x$'s coordinates, which only helps), producing at least $\mu$-many values
of $y$ in $c_x$. Concretely: for $y=(h_1,h_2)$, $x=(h_3,h_4)$, take
$h_3$ with $\mu$ predecessors in $H$ and count the pairs from
$H\cap h_3$; for $y=(h_1,h_3)$, $x=(h_2,h_4)$, take $h_2$ with $\mu$
predecessors and $h_4$ large, and count $h_1\in H\cap h_2$ (with any fixed
admissible $h_3$); for $y=(h_2,h_3)$, $x=(h_1,h_4)$, take $h_4$ so that
$|H\cap(h_1,h_4)|\ge\mu$ and count the pairs between.

\emph{Three-element patterns}: if $y$'s non-shared coordinate can be
varied upward (patterns $y=(h_1,h_3)$ or $(h_2,h_3)$ with $h_3$ free
above $x$, and the shared-coordinate analogues), fix $x$ and vary it:
$\kappa$-many $y$'s. If $y$'s non-shared coordinate lies below a
coordinate of $x$ (patterns $y=(h_1,h_2)$ or $(h_1,h_3)$ with $h_1$ free
below), choose the relevant coordinate of $x$ with $\mu$ predecessors in
$H$ and count.

Every ``yes'' is thus impossible; all patterns are homogeneously
``$y\notin c_x$,'' and $S$ mutually avoids the assignment.
\end{proof}

\begin{remark}
The proof uses nothing about $\D\otimes\D$ beyond: it contains the upper
triangles of Rowbottom-homogeneous sets. It therefore applies verbatim to
all finite Fubini powers $\D^{\otimes n}$ (color $2n$-element sets by
membership patterns; every ``yes'' pattern forces a large $c_x$ by
freezing $x$ and freeing a coordinate of $y$, upward or below a
high pivot), and to any ultrafilter extending the filter of triangle-sets
of $\D$-homogeneous sets. We will not need this generality.
\end{remark}

\section{Transfer, and the answer to Question 5.5 of \cite{FWIII}}
\label{sec:transfer}

\begin{definition}\label{def:EPP}
A $\kappa$-complete nonprincipal ultrafilter $\U$ on a set $X$ of size
$\kappa$ has the \emph{exact-paving property} (EPP) if for every
assignment of countable sets $c_x\subseteq X$ to points $x\in X$ there is
$S\in\U$ with $y\notin c_x$ for all distinct $x,y\in S$.
\end{definition}

Thus \cite{FWIII}, Theorem 2.1 says normal measures have EPP (via the
diagonal intersection), and Theorem~\ref{thm:paving-product} says
$\D\otimes\D$ has EPP (via homogeneity). The point of the definition is
that EPP is exactly what the positive theory of \cite{FWIII} consumes:

\begin{theorem}\label{thm:main}
Let $\U$ be a $\kappa$-complete nonprincipal ultrafilter with EPP on an
index set $X$, $|X|=\kappa$, and let $\Phi(\U)$ be the block filter in
$\PP(\ell^2(X))$. Then every conclusion of \cite{FWIII}, Sections 2--3
holds for $\U$: exact paving of every bounded operator on a measure-one
block; the exact dichotomy for closed subspaces; the diagonal-flattening
estimate; and $\Phi(\U)$ is a proper, upward closed,
${<}\kappa$-intersection-closed, nonprincipal filter deciding every closed
subspace --- a $\kappa$-complete nonprincipal pluf --- with
$S_{\Phi(\U)}$ a singleton and the round-slice property for every
operator. In particular, all of this holds for $\U=\D\otimes\D$.
\end{theorem}

\begin{proof}
Exact paving of a bounded operator $T$ asks precisely for a mutually
avoiding set for the assignment
$c_x=\supp(Te_x)\cup\supp(T^*e_x)$, which is countable because vectors of
$\ell^2(X)$ have countable support; this is EPP verbatim. Every subsequent
argument of \cite{FWIII}, Sections 2--3 --- the orthogonality of the
defect vectors $P_{W^\perp}e_x$ over the paving set, the exact
computation $W\cap\HH_{S_0}=\HH_{S_1}$, the ultrafilter decision, the
filter axioms (which use $\kappa$-completeness and, for nonprincipality,
that countable sets are $\U$-small, immediate from $\kappa$-completeness
and nonprincipality), the flattening (which uses paving, the existence of
ultrafilter limits of bounded functions, and a termwise series
comparison), and the singleton face (paving plus flattening plus the
state estimates of \cite{FWI}) --- consumes normality only through the
paving conclusion, as one verifies by inspection; indeed the mechanized
development accompanying \cite{FWIII} already isolates the dependence,
deriving the dichotomy, the flattening and the filter packaging from the
paving statement together with $\kappa$-completeness-type hypotheses
alone. Theorem~\ref{thm:paving-product} supplies paving for
$\D\otimes\D$, and $\D\otimes\D$ is $\kappa$-complete and nonprincipal.
\end{proof}

\begin{corollary}[Answer to {\cite{FWIII}}, Question 5.5]
\label{cor:answer}
Maximality of the block filter does not characterize normality up to
isomorphism: $\Phi(\D\otimes\D)$ is a $\kappa$-complete nonprincipal pluf
with singleton face, while $\D\otimes\D$ is isomorphic to no normal
measure (Proposition~\ref{prop:notiso}).
\end{corollary}

The invariants now sit as follows, all implications for
$\kappa$-complete nonprincipal ultrafilters on a set of size $\kappa$:
\[
\text{full-selector}\;\Longleftarrow\;\text{normal}
\;\Longrightarrow\;\text{EPP}
\;\Longrightarrow\;\text{$\Phi(\U)$ maximal}
\;\Longrightarrow\;\text{$\sigma$-Q},
\]
where the last implication is \cite{FWIII}, Proposition 5.3, and EPP
implies $\sigma$-Q directly: given a countable-piece partition, apply EPP
to $c_x=\sigma(x)\setminus\{x\}$; a mutually avoiding set is a partial
selector. The product has EPP, maximality and $\sigma$-Q while failing
the full-selector property; so maximality detects strictly less than
normality. Whether the two middle implications reverse we do not know
(Question~\ref{q:EPP}).

\begin{remark}[The two cardinals compared]\label{rem:compare}
At $\aleph_0$, the block filter of an ultrafilter (completed by the
finite-codimension subspaces) is maximal exactly for selective
ultrafilters (\cite{FWII}, Theorem 3.4), and products $U\otimes V$ are
never selective; at $\kappa$ along $\kappa$-complete ultrafilters,
products of a normal measure \emph{are} maximal. The linear theory at
measurable cardinals thus sees Rudin--Keisler structure more coarsely
than the classical theory at $\omega$: the class of ``tame''
ultrafilters strictly exceeds the RK-minimal ones. The mechanism is
plain in the proofs: at $\omega$, maximality needs the full
Ramsey-theoretic strength of selectivity through Mathias's theorem,
while at $\kappa$ the countable-support phenomenon reduces maximality to
paving, and paving to homogeneity --- which, unlike normality, survives
finite products.
\end{remark}

\section{Constraints on blockers at $\aleph_0$}\label{sec:rigidity}

We now return to the rigidity conjecture of \cite{FWII}. Fix a
non-selective ultrafilter $\U$ on $\N$, a witness partition
$\N=\bigsqcup_kA_k$ (no piece and no selector in $\U$), the constraint
vectors $f_k=\sum_{n\in A_k}2^{-n}e_n$, and $W=\{f_k:k\}^\perp$; recall
(\cite{FWII}, Theorem 3.4) that $W$ is addable to the block filter
$\Phi(\U)$ but not a member. Call a closed subspace $R$ an \emph{addable
blocker} if $R\cap W=0$ and $\dim(R\cap\HH_S)=\infty$ for every
$S\in\U$; the rigidity conjecture (\cite{FWII}, Conjecture 3.8) asserts
none exists, equivalently that $W$ belongs to every pluf extension of
$\Phi(\U)$. The results of \cite{FWII} exclude blockers spanned by
disjointly supported families and blockers meeting a single piece-block
two-dimensionally. Here are three further constraints. For
$T\subseteq\N$ write $\HH_T$ for the block, and for $x\in\HH$ let the
\emph{piece-spread} of $x$ be $\{k:x|_{A_k}\neq0\}$.

\begin{lemma}[Generalized thinness]\label{lem:genthin}
If $R$ is closed with $R\cap W=0$, then for every finite set $F$ of
piece indices,
\[
\dim\Bigl(R\cap\HH_{\bigcup_{k\in F}A_k}\Bigr)\;\le\;|F|.
\]
\end{lemma}

\begin{proof}
The $|F|$ linear functionals $\langle\,\cdot\,,f_k\rangle$, $k\in F$,
cut any subspace of $\HH_{\cup F}$ of dimension exceeding $|F|$ in a
nonzero vector $v$ orthogonal to all $f_k$, $k\in F$; and $v\perp f_j$
for $j\notin F$ automatically (disjoint supports), so
$v\in R\cap W=0$.
\end{proof}

Taking $F$ a singleton recovers \cite{FWII}, Lemma 3.6.

\begin{lemma}[Support]\label{lem:support}
If $R$ is an addable blocker, then
$T^*:=\bigcup_{x\in R}\supp x\in\U$. Consequently no addable blocker is
contained in $\HH_T$ for any $T\notin\U$ --- in particular, none is
contained in the block of a partial selector, or of any union of a
partial selector with a set outside $\U$.
\end{lemma}

\begin{proof}
If $T^*\notin\U$ then $S_0=\N\setminus T^*\in\U$; but
$R\subseteq\HH_{T^*}$, so a nonzero $x\in R\cap\HH_{S}$ for
$S\in\U$ has infinite support inside $T^*\cap S$. More precisely, for
any $T\supseteq T^*$ with $T\notin\U$: for every $S\in\U$,
$R\cap\HH_S\subseteq\HH_{T\cap S}$, and $\dim(R\cap\HH_S)=\infty$
forces $T\cap S$ infinite; ``$T\cap S$ infinite for all $S\in\U$'' is
equivalent, for an ultrafilter, to $T\in\U$ (otherwise
$\N\setminus T\in\U$ and $T\cap(\N\setminus T)=\emptyset$). So
$T^*\in\U$.
\end{proof}

\begin{lemma}[Piece-spread]\label{lem:spread}
If $R$ is an addable blocker, then for every $S\in\U$ the space
$R\cap\HH_S$ contains a vector of infinite piece-spread.
\end{lemma}

\begin{proof}
Fix $S\in\U$ and write $R_S=R\cap\HH_S$, a closed, hence completely
metrizable, subspace with $\dim R_S=\infty$. The vectors of $R_S$ of
piece-spread contained in a fixed finite $F$ form the closed set
$E_F=R_S\cap\HH_{\bigcup_{k\in F}A_k}$, of dimension at most $|F|$ by
Lemma~\ref{lem:genthin}; a finite-dimensional proper closed subspace of
an infinite-dimensional Banach space is nowhere dense. The vectors of
finite piece-spread form the countable union $\bigcup_F E_F$ over
finite $F$, which by the Baire category theorem is not all of $R_S$.
\end{proof}

\begin{remark}\label{rem:shape}
Together with \cite{FWII}, Proposition 3.7, the lemmas confine any
counterexample to the conjecture to a narrow shape: a blocker must have
total support in $\U$ (Lemma~\ref{lem:support}), meet every finite union
of piece-blocks in dimension bounded by the number of pieces
(Lemma~\ref{lem:genthin}), be spanned by no disjointly supported family
(\cite{FWII}), and sustain, inside every measure-one block
simultaneously, infinitely many dimensions of vectors each spread over
infinitely many pieces (Lemma~\ref{lem:spread}) --- that is, its
intersection with every $\HH_S$ must be produced by cancellation among
overlapping supports across infinitely many pieces, never by whole
generators. The measurable-cardinal results of
Sections~\ref{sec:inj}--\ref{sec:transfer} bound the conjecture from above:
along $\D$ or $\D\otimes\D$ the $\sigma$-Q-point property holds, so no
witness subspace exists and the rigidity question does not arise.
Rigidity, if it holds, is a phenomenon of small cardinality, tied to the
same Ramsey-theoretic poverty of $\omega$ that makes selectivity a
special property there and a theorem at $\kappa$.
\end{remark}

\section{Questions}\label{sec:questions}

\begin{question}\label{q:EPP}
For $\kappa$-complete nonprincipal ultrafilters on a set of size
$\kappa$: does maximality of the block filter imply EPP? Does the
$\sigma$-Q-point property imply EPP (equivalently, given
\cite{FWIII}, Proposition 5.3 and Theorem~\ref{thm:main}, are
$\sigma$-Q, EPP and maximality all equivalent)? A negative answer to
either would require an ultrafilter separating selector-type from
free-set-type combinatorics at $\kappa$; candidates presumably live
among the many $\kappa$-complete ultrafilters of a Kunen-style model
rather than among iterated products, which Theorem
\ref{thm:paving-product} and its remark place on the EPP side.
\end{question}

\begin{question}\label{q:mixed}
Do mixed products $\D_1\otimes\D_2$ of two distinct normal measures on
the same $\kappa$ have EPP? The homogeneity argument as given requires
one set $H$ serving both coordinates, i.e.\ Rowbottom-homogeneous sets
in $\D_1\cap\D_2$, which need not exist. A proof or a counterexample
would locate how much of Theorem~\ref{thm:paving-product} is about
products and how much about self-products.
\end{question}

\begin{question}\label{q:rigidity}
The rigidity conjecture of \cite{FWII} stands: does some non-selective
$\U$ admit an addable blocker of its witness subspace --- necessarily of
the shape delimited in Remark~\ref{rem:shape} --- or does $W$ belong to
every maximalization of every block filter of a non-selective
ultrafilter? A first target consistent with all known constraints: a
blocker for a product ultrafilter $\U=V\otimes V$, built from vectors
spread along the graphs of finitely-branching trees of pieces, where the
selector obstructions of \cite{FWII}, Proposition 3.7 do not directly
apply.
\end{question}

\subsection*{Formal verification}

Every theorem of this paper has been formalized and machine-checked in
Lean~4 with Mathlib (toolchain v4.28.0), extending the artifacts of
\cite{FWIII}: Lemma~\ref{lem:full-selector}, without countability or
nonemptiness of the pieces; Proposition~\ref{prop:notiso}, through an
isomorphism-invariance lemma for the full-selector property;
Theorem~\ref{thm:inj} and Corollary~\ref{cor:product-sigmaQ};
Theorem~\ref{thm:paving-product}, as the enumeration of twelve ordered
placement patterns, of which four are the mid-region patterns treated by the
tail-then-pivot step of Section~\ref{sec:paving}; the exact-paving property
with the implication EPP $\Rightarrow$ $\sigma$-Q; the lattice-level clauses
of Theorem~\ref{thm:main} --- operator paving, the exact dichotomy, the
flattening estimate, and the decision property and nonprincipality of the
block filter, instantiated at the Fubini product --- and
Lemmas~\ref{lem:genthin}--\ref{lem:spread}. Classical inputs are
quarantined as hypotheses (Rowbottom homogeneity for triples and
quadruples, membership of tails, uncountable pivots, and the Fodor
property), so every mechanized statement is a theorem of ZFC. The
injectivity theorem consumes the tails hypothesis, which holds for every
uniform ultrafilter and is invisible at a normal measure. The artifacts
compile with no unproved obligations, and every statement's axiom footprint
is exactly \texttt{propext}, \texttt{Classical.choice},
\texttt{Quot.sound}. Two items are not covered: the finite-power remark of
Section~\ref{sec:paving}, and the state-theoretic clauses of
Theorem~\ref{thm:main} (singleton face, round slices). The face machinery
those clauses invoke is itself verified in \cite{FWI}; its transfer to
ultrafilters given only by EPP is not part of the present development.

\subsection*{Acknowledgements}

Portions of these results were developed in the course of an extended
dialogue with Anthropic's Claude language model; responsibility for
correctness rests with the authors.

\begin{thebibliography}{9}

\bibitem{FWI} D.~V.~Feldman and A.~Wilce, \emph{Ultrafilter limits in
the projection lattice of a Hilbert space}, preprint (paper I of this
series).

\bibitem{FWII} D.~V.~Feldman and A.~Wilce, \emph{Intimate subspaces,
block filters, and the diagonalization of maximal projection filters},
preprint (paper II of this series).

\bibitem{FWIII} D.~V.~Feldman and A.~Wilce, \emph{Projection-lattice
ultrafilters at a measurable cardinal}, preprint (paper III of this
series).

\bibitem{plufLean} D.~V.~Feldman and A.~Wilce, \emph{pluf-lean: a Lean~4 verification of maximal filters in the projection lattice of a Hilbert space}, Zenodo, \href{https://doi.org/10.5281/zenodo.21987888}{\texttt{10.5281/zenodo.21987888}}.

\bibitem{Kanamori} A.~Kanamori, \emph{The Higher Infinite}, 2nd ed.,
Springer Monographs in Mathematics, Springer, 2003.

\bibitem{Rowbottom} F.~Rowbottom, \emph{Some strong axioms of infinity
incompatible with the axiom of constructibility}, Ann. Math. Logic
\textbf{3} (1971), 1--44.

\end{thebibliography}

\end{document}
